\chapter{Units}
\begin{center}\itshape
The smallest solution of $x^2-61y^2=1$ in positive integers is $x=1\,766\,319\,049$, $y=226\,153\,980$.\\
Why must there be one at all, and what replaces Pell's equation in a field of degree $3$ or $4$?
\end{center}

In February 1657 Pierre de Fermat sent a challenge to the mathematicians of Europe, and in particular, with some malice,
to those of England: find a square which, multiplied by $61$ and increased by $1$, gives a square. He added, to spare
them too much effort, that $109$ is another good example. He knew exactly what he was doing: the smallest answers have
ten and fifteen digits. Brouncker and Wallis eventually found a general method, and Euler later attached the equation,
by mistake, to the name of John Pell, who had nothing to do with it.

None of them knew that the problem had been solved in India five centuries earlier. Brahmagupta, in 628, had said that
whoever solves $x^2-92y^2=1$ within a year is a true mathematician; and in 1150 Bh\=askara~II solved $x^2-61y^2=1$
itself with his \emph{chakravala}, the ``cyclic method'', finding $x=1\,766\,319\,049$. Fermat's trap had been sprung long
before he set it.

In the language of this volume, solutions of Pell's equation are \emph{units}: $x+y\sqrt{61}$ is invertible in $\Z[\sqrt{61}]$
exactly when $x^2-61y^2=\pm1$. Chapter~\ref{ch:quadratic} proves that the units of a real quadratic field are $\pm\varepsilon^k$ for one
fundamental unit $\varepsilon$ (Theorem~\ref{thm:realunits}), found by continued fractions. What happens in degree $3$, or $10$?
In 1846 Dirichlet found the answer, according to the legend while listening to the Easter music in the Sistine Chapel: the
unit group is a finite group of roots of unity times a free abelian group of rank $r_1+r_2-1$, the number of
``infinite places'' minus one.

The proof is the second great application of Minkowski's theorem. Taking logarithms of absolute values turns the
multiplicative group of units into an additive group of vectors: it lies in a hyperplane, it is discrete, and Minkowski
provides enough units to fill the hyperplane. The volume of the resulting lattice, the \emph{regulator}, is the second
invariant of a number field after the discriminant, and it sits next to the class number in the class number formula of
Chapter~11.

\begin{reminder}[from Volumes 1 and 2, and Chapters 3--5]
\emph{Units} (Chapter~3, Proposition~\ref{prop:units}): $\alpha\in\mathcal O_K$ is a unit iff $N(\alpha)=\pm1$. A finite
subgroup of the multiplicative group of a field is cyclic (Volume~2, Theorem~9.4). A subgroup of a free abelian group of
rank $m$ is free of rank $\le m$ (Volume~2, Lemma~7.5). \emph{Chapter~5}: the Minkowski embedding $j$, the covolume
$2^{-r_2}\sqrt{|d_K|}$ of $j(\mathcal O_K)$, and the convex body theorem. \emph{Chapter~4}: an integer $m\ge1$ lies in only
finitely many ideals; $N(\alpha\mathcal O)=|N(\alpha)|$.
\end{reminder}

\section{Roots of unity and the logarithmic embedding}

Let $K$ have signature $(r_1,r_2)$, $n=r_1+2r_2$, real embeddings $\sigma_1,\dots,\sigma_{r_1}$ and one embedding $\sigma_{r_1+j}=\tau_j$ in
each complex pair. Put $r=r_1+r_2-1$ and $e_i=1$ for $i\le r_1$, $e_i=2$ for $i>r_1$, so that $|N(x)|=\prod_{i=1}^{r+1}
|\sigma_i(x)|^{e_i}$. Write $\mu_K$ for the group of roots of unity in $K$.

\begin{definition}
The \emph{logarithmic embedding} is the group homomorphism
\[
\ell:\mathcal O_K^\times\to\R^{r+1},\qquad u\mapsto\bigl(e_1\log|\sigma_1(u)|,\ \dots,\ e_{r+1}\log|\sigma_{r+1}(u)|\bigr).
\]
\end{definition}

It is a homomorphism because $|\sigma_i(uv)|=|\sigma_i(u)||\sigma_i(v)|$. Its image lies in the hyperplane
$H=\{x\in\R^{r+1}:\sum x_i=0\}$, since the sum of the coordinates is $\log|N(u)|=\log1=0$.

\begin{proposition}\label{prop:logkernel}
\begin{enumerate}[label=(\alph*)]
\item For every $M>0$ there are only finitely many $u\in\mathcal O_K^\times$ with $\|\ell(u)\|_\infty\le M$. In particular
$\ell(\mathcal O_K^\times)$ is a discrete subgroup of $H$.
\item The kernel of $\ell$ is $\mu_K$, which is finite and cyclic.
\end{enumerate}
\end{proposition}
\begin{proof}
(a) If $\|\ell(u)\|_\infty\le M$, all conjugates of $u$ have absolute value at most $e^M$, so the coefficients of its
characteristic polynomial $\chi_u=\prod(X-\sigma_i(u))$, integers by Proposition~\ref{prop:minpolyZ}, are bounded by
$\binom nke^{kM}$. There are finitely many such polynomials, and each has at most $n$ roots.

(b) A root of unity $\zeta$ has $|\sigma_i(\zeta)|=1$ for all $i$, so $\ell(\zeta)=0$. Conversely the kernel is finite by
(a) with $M=0$; it is a finite group, so each of its elements has finite order, and it consists of roots of unity. A
finite subgroup of $K^\times$ is cyclic (Volume~2, Theorem~9.4).
\end{proof}

\begin{lemma}\label{lem:discrete}
A discrete subgroup $\Gamma$ of $\R^m$ is free abelian, generated by linearly independent vectors.
\end{lemma}
\begin{proof}
Let $V$ be the span of $\Gamma$ and $v_1,\dots,v_s\in\Gamma$ a basis of $V$. Let $\Gamma_0=\bigoplus\Z v_i\subseteq\Gamma$
and $F=\{\sum t_iv_i:0\le t_i<1\}$. Every $\gamma\in\Gamma$ is $\gamma_0+f$ with $\gamma_0\in\Gamma_0$ and $f\in F\cap\Gamma$
(split the coordinates into integer and fractional parts). As $F$ is bounded and $\Gamma$ discrete, $F\cap\Gamma$ is finite, so
$N=[\Gamma:\Gamma_0]$ is finite and $N\Gamma\subseteq\Gamma_0$. Then $\Gamma\subseteq\frac1N\Gamma_0\cong\Z^s$ is free of rank
$\le s$ (Volume~2, Lemma~7.5), and its rank is $s$ since it contains $\Gamma_0$. A $\Z$-basis of $\Gamma$ spans $V$ and has
$s=\dim V$ elements, so it is linearly independent.
\end{proof}

So $\ell(\mathcal O_K^\times)$ is free of rank at most $\dim H=r$. The whole difficulty is to show that it is \emph{exactly}
$r$: we need many units.

\section{Dirichlet's unit theorem}

\begin{lemma}\label{lem:smallbox}
Let $B=(\frac2\pi)^{r_2}\sqrt{|d_K|}$. If $c_1,\dots,c_{r+1}>0$ satisfy $\prod_ic_i^{e_i}\ge B$, there is $\alpha\in\mathcal
O_K\setminus\{0\}$ with $|\sigma_i(\alpha)|\le c_i$ for all $i$, and then $|N(\alpha)|\le\prod c_i^{e_i}$.
\end{lemma}
\begin{proof}
In the Minkowski embedding, the set $\{|y_i|\le c_i\ (i\le r_1),\ |z_j|\le c_{r_1+j}\}$ is compact, convex, symmetric, of
volume $\prod_{i\le r_1}2c_i\prod_j\pi c_{r_1+j}^2=2^{r_1}\pi^{r_2}\prod c_i^{e_i}$. This is at least $2^n\cdot2^{-r_2}\sqrt{|d_K|}=
2^n\operatorname{covol}(j(\mathcal O_K))$ exactly when $\prod c_i^{e_i}\ge2^{n-r_1-r_2}\pi^{-r_2}\sqrt{|d_K|}=B$ (as
$n-r_1-r_2=r_2$). Apply Theorem~\ref{thm:minkowski}(b).
\end{proof}

\begin{lemma}\label{lem:onebig}
For each $k\in\{1,\dots,r+1\}$ there is a unit $u$ with $|\sigma_i(u)|<1$ for every $i\ne k$.
\end{lemma}
\begin{proof}
We build a sequence $\alpha_1,\alpha_2,\dots$ of non-zero integers with $|N(\alpha_m)|\le B$ and $|\sigma_i(\alpha_{m+1})|<
|\sigma_i(\alpha_m)|$ for all $i\ne k$. Start with any $\alpha_1$ given by Lemma~\ref{lem:smallbox}. Given $\alpha_m$, put
$c_i=\frac12|\sigma_i(\alpha_m)|$ for $i\ne k$, and choose $c_k$ so that $\prod c_i^{e_i}=B$; Lemma~\ref{lem:smallbox} gives
$\alpha_{m+1}$ with $|\sigma_i(\alpha_{m+1})|\le c_i<|\sigma_i(\alpha_m)|$ for $i\ne k$ and $|N(\alpha_{m+1})|\le B$.

The ideals $\alpha_m\mathcal O$ have norm $|N(\alpha_m)|\le B$, and there are finitely many ideals of norm at most $B$
(each contains its norm). So $\alpha_m\mathcal O=\alpha_{m'}\mathcal O$ for some $m<m'$, and $u=\alpha_{m'}/\alpha_m$ is a
unit. For $i\ne k$, $|\sigma_i(u)|=|\sigma_i(\alpha_{m'})|/|\sigma_i(\alpha_m)|<1$.
\end{proof}

\begin{lemma}\label{lem:dominant}
Let $A=(a_{ij})$ be a real $r\times r$ matrix with $a_{ij}<0$ for $i\ne j$ and $\sum_ja_{ij}>0$ for every $i$. Then $A$
is invertible.
\end{lemma}
\begin{proof}
Suppose $Ax=0$ with $x\ne0$; replacing $x$ by $-x$, let $x_i>0$ be a coordinate of largest absolute value. Then
$0=\sum_ja_{ij}x_j\ge a_{ii}x_i+\sum_{j\ne i}a_{ij}x_i=x_i\sum_ja_{ij}>0$, using $a_{ij}x_j\ge a_{ij}x_i$ for $j\ne i$
(as $a_{ij}<0$ and $x_j\le x_i$). A contradiction.
\end{proof}

\begin{theorem}[Dirichlet's unit theorem]\label{thm:dirichlet}
$\ell(\mathcal O_K^\times)$ is a lattice in the hyperplane $H$, and
\[
\mathcal O_K^\times\cong\mu_K\times\Z^{r},\qquad r=r_1+r_2-1.
\]
That is, there are \emph{fundamental units} $\varepsilon_1,\dots,\varepsilon_r$ such that every unit is uniquely
$\zeta\varepsilon_1^{k_1}\cdots\varepsilon_r^{k_r}$ with $\zeta\in\mu_K$ and $k_i\in\Z$.
\end{theorem}
\begin{proof}
For $k=1,\dots,r$ take $u_k$ from Lemma~\ref{lem:onebig}. The vector $\ell(u_k)$ has negative coordinates except the $k$-th,
and its coordinates sum to $0$. Let $A$ be the $r\times r$ matrix of the first $r$ coordinates of $\ell(u_1),\dots,
\ell(u_r)$ (row $k$ for $u_k$). Off the diagonal its entries are negative, and each row sum equals minus the omitted
$(r+1)$-th coordinate, which is positive. By Lemma~\ref{lem:dominant}, $A$ is invertible, so the $\ell(u_k)$ are linearly
independent, and $\ell(\mathcal O_K^\times)$ spans $H$. By Proposition~\ref{prop:logkernel} and Lemma~\ref{lem:discrete},
it is free with a basis $\ell(\varepsilon_1),\dots,\ell(\varepsilon_r)$ of $H$: a lattice in $H$.

Every unit $u$ then has $\ell(u)=\sum k_i\ell(\varepsilon_i)$ for unique integers $k_i$, and $u\varepsilon_1^{-k_1}\cdots
\varepsilon_r^{-k_r}$ is in $\ker\ell=\mu_K$. Uniqueness of $\zeta$ and the $k_i$ follows from the uniqueness of the $k_i$.
\end{proof}

\begin{think}
For which fields is the unit group finite? List them: they are $\Q$ and two other kinds. Then show that a totally real
field of degree $3$ has two independent units, and a field of degree $4$ with signature $(2,1)$ also has two.
\end{think}

\begin{figure}[ht]
\centering
\fig{ch04-loglattice}
\caption{The units of the real cubic field $K=\Q(\theta)$, $\theta=2\cos\frac{2\pi}7$, seen through the logarithmic
embedding: the points $\ell(\pm\theta^a(1+\theta)^b)$, drawn in the plane $x_1+x_2+x_3=0$ of $\R^3$. They form a lattice,
generated by $\ell(\theta)$ and $\ell(1+\theta)$; the shaded parallelogram has area $\sqrt3\,R_K\approx0.910$, with
$R_K\approx0.5255$ the regulator (Exercise~\ref{ex:cyclic7}).}
\label{fig:loglattice}
\end{figure}

\section{The regulator}

\begin{definition}
Let $\varepsilon_1,\dots,\varepsilon_r$ be fundamental units. The \emph{regulator} $R_K$ is the absolute value of the
determinant of the $r\times r$ matrix obtained from the $r\times(r+1)$ matrix $\bigl(e_i\log|\sigma_i(\varepsilon_k)|
\bigr)$ by deleting one column. For $r=0$, $R_K=1$.
\end{definition}

\begin{proposition}\label{prop:regwelldef}
$R_K$ depends neither on the choice of fundamental units nor on the deleted column, and
$\operatorname{covol}_H(\ell(\mathcal O_K^\times))=\sqrt{r+1}\,R_K$, the covolume being measured in the hyperplane $H$.
\end{proposition}
\begin{proof}
Two systems of fundamental units differ, modulo $\mu_K$, by a matrix of $\mathrm{GL}_r(\Z)$, which multiplies every
$r\times r$ minor by $\pm1$. Let $L$ be the $r\times(r+1)$ matrix of the $\ell(\varepsilon_k)$, and add to it the row
$(1,\dots,1)$, getting a square matrix $\tilde L$. Since the rows of $L$ have sum $0$, adding all columns of $\tilde L$ to
the $c$-th column replaces it by $(0,\dots,0,r+1)^{\mathsf T}$; expanding along it, $\det\tilde L=\pm(r+1)\det L_c$, where
$L_c$ is $L$ with column $c$ deleted. So all the $|\det L_c|$ are equal to $|\det\tilde L|/(r+1)$. Finally, the vector
$\frac1{\sqrt{r+1}}(1,\dots,1)$ is a unit normal to $H$, so the $(r+1)$-dimensional volume spanned by it and the
$\ell(\varepsilon_k)$, namely $|\det\tilde L|/\sqrt{r+1}$, is the $r$-dimensional covolume of $\ell(\mathcal O_K^\times)$ in
$H$. That is $\sqrt{r+1}\,R_K$.
\end{proof}

For a real quadratic field with fundamental unit $\varepsilon>1$, $R_K=\log\varepsilon$. For $\Q(\sqrt2)$, $R_K=\log(1+\sqrt2)
\approx0.881$; for $\Q(\sqrt{61})$, the fundamental unit is $\varepsilon=\frac{39+5\sqrt{61}}2$, of norm $-1$, and
$R_K=\log\varepsilon\approx3.664$. Exercise~\ref{ex:pell61} explains how $\varepsilon^6=1\,766\,319\,049+226\,153\,980\sqrt{61}$ solves the
opening question.

\section{Finding fundamental units}

Dirichlet's proof is not constructive in any useful way. For real quadratic fields, Volume~1 found $\varepsilon$ through
the continued fraction of $\sqrt d$. This also gives an efficient algorithm. For other fields one finds \emph{some} units,
by searching for elements of small norm and dividing them as in Lemma~\ref{lem:onebig}, and then one proves they are
fundamental with a lower bound for the regulator. Here is the classical bound for cubic fields with one real embedding,
where $r=1$.

\begin{proposition}[an Artin-type bound]\label{prop:artin}
Let $K$ be a cubic field with $r_1=r_2=1$, and $\varepsilon>1$ a unit of $K$ (in its real embedding). Then
$|d_K|<4\varepsilon^3+35$.
\end{proposition}
\begin{proof}
As $\varepsilon$ is irrational, $K=\Q(\varepsilon)$. Its other conjugates are $\rho e^{\pm i\varphi}$ with $\rho>0$ and
$\sin\varphi\ne0$, and $N(\varepsilon)=\varepsilon\rho^2>0$, so $N(\varepsilon)=1$ and $\rho=\varepsilon^{-1/2}$. By
Proposition~\ref{prop:disc}(c),
\[
\operatorname{disc}(1,\varepsilon,\varepsilon^2)=\bigl[(\varepsilon-\rho e^{i\varphi})(\varepsilon-\rho e^{-i\varphi})
(2i\rho\sin\varphi)\bigr]^2.
\]
Now $(\varepsilon-\rho e^{i\varphi})(\varepsilon-\rho e^{-i\varphi})=\varepsilon^2-2\varepsilon\rho\cos\varphi+\rho^2=
\rho^{-1}(x+x^{-1}-2\cos\varphi)$ with $x=\rho^{-3}=\varepsilon^{3/2}>1$. So, with $s=\sin\varphi$, $c=\cos\varphi$,
\[
\bigl|\operatorname{disc}(1,\varepsilon,\varepsilon^2)\bigr|=4E,\qquad E=s^2(x+x^{-1}-2c)^2.
\]
Expanding and writing $s^2x^2=x^2-c^2x^2$,
\[
E=x^2-c^2x^2-4cs^2x-4cs^2x^{-1}+s^2(2+x^{-2}+4c^2).
\]
Complete the square: $-c^2x^2-4cs^2x=4s^4-(cx+2s^2)^2\le4s^4$. With $x>1$, $-4cs^2x^{-1}\le4s^2|c|$ and $x^{-2}<1$.
Hence $E<x^2+4s^4+4s^2|c|+3s^2+4s^2c^2=x^2+7s^2+4s^2|c|$, using $s^2+c^2=1$. Finally $s^2\le1$ and $s^2|c|=(1-c^2)|c|
\le\frac2{3\sqrt3}$ (maximum at $c^2=\frac13$), so $E<x^2+7+\frac8{3\sqrt3}<x^2+8.6$. Therefore $|\operatorname{disc}(1,
\varepsilon,\varepsilon^2)|<4\varepsilon^3+34.4$, and $|d_K|\le|\operatorname{disc}(1,\varepsilon,\varepsilon^2)|$ by the index
formula.
\end{proof}

Artin's sharper version has $24$ in place of $35$ (Cassels and Fr\"ohlich, or Marcus, Chapter~5), but ours is enough here.

\begin{example}[the units of {$\Q(\sqrt[3]2)$}]\label{exa:unitscube}
Here $r=1$ and $\mu_K=\{\pm1\}$ ($K\subseteq\R$). The element $\varepsilon=1+\sqrt[3]2+\sqrt[3]4$ satisfies
$\varepsilon(\sqrt[3]2-1)=2-1=1$, so it is a unit, $\varepsilon\approx3.847$. Let $\eta>1$ be the fundamental unit, so
$\varepsilon=\eta^k$ with $k\ge1$. By Proposition~\ref{prop:artin}, $4\eta^3>108-35=73$, so $\eta>2.63$ and $\eta^2>6.9>
\varepsilon$. Hence $k=1$: the units of $\Z[\sqrt[3]2]$ are exactly $\pm(1+\sqrt[3]2+\sqrt[3]4)^k$, $k\in\Z$, and
$R_K=\log\varepsilon\approx1.3474$. Exercise~\ref{ex:normform} of Chapter~3 found infinitely many solutions of the cubic
Pell equation $a^3+2b^3+4c^3-6abc=1$; now we know them all: they are the coefficients of the powers of $\varepsilon$.
\end{example}

\section{Exercises}

\begin{exercise}[$\star$]\label{ex:muK}
Show that if $K$ contains a primitive $m$-th root of unity then $\varphi(m)\mid n$. Deduce that $\mu_K=\{\pm1\}$ for every
quadratic field except $\Q(i)$ and $\Q(\sqrt{-3})$, and for every field of odd degree.
\end{exercise}

\begin{exercise}[$\star$]\label{ex:pell61}
Check that $\varepsilon=\frac{39+5\sqrt{61}}2$ is a unit of norm $-1$. Show that for $d\equiv1\pmod4$, the cube of every unit
of $\mathcal O_{\Q(\sqrt d)}$ lies in $\Z[\sqrt d]$. Deduce that the smallest solution of $x^2-61y^2=1$ with $y>0$ comes
from $\varepsilon^6$, and compute it.
\end{exercise}

\begin{exercise}[$\star\star$]\label{ex:cyclic7}
Let $\theta=2\cos\frac{2\pi}7$, a root of $X^3+X^2-2X-1$, and $K=\Q(\theta)$. Show that $\theta$ and $1+\theta$ are units, that
one other root is $\theta^2-2$, find the third, and compute numerically the determinant giving the
regulator of the subgroup they generate. Show that they are independent.
\end{exercise}

\begin{exercise}[$\star\star$]\label{ex:regcol}
Check Proposition~\ref{prop:regwelldef} on $\Q(\sqrt2)$: with $\varepsilon=1+\sqrt2$, compute $\ell(\varepsilon)$, both
$1\times1$ minors, and the covolume of $\ell(\mathcal O^\times)$ on the line $x_1+x_2=0$.
\end{exercise}

\begin{exercise}[$\star\star$]\label{ex:sunits}
\emph{$S$-units.} Let $S$ be a finite set of primes of $\mathcal O_K$, and $\mathcal O_S^\times$ the group of $x\in K^\times$ with
$v_P(x)=0$ for all $P\notin S$. Show that $\mathcal O_S^\times\cong\mu_K\times\Z^{r+|S|}$. \emph{(Map $x$ to $(v_P(x))_{P\in S}$
and use the finiteness of the class group.)}
\end{exercise}

\begin{exercise}[$\star\star$]\label{ex:thue}
\emph{Thue's equation $x^3-2y^3=1$.} Show that the integer solutions correspond to the integers $k$ for which
$\varepsilon^{k}$ (with $\varepsilon$ as in Example~\ref{exa:unitscube}) has the form $x-y\sqrt[3]2$, with no $\sqrt[3]4$
term. Find the two values of $k$ with $|k|\le3$ that work, and the corresponding solutions. (Skolem's $p$-adic method shows that there are no others.)
\end{exercise}

\begin{exercise}[$\star\star\star$]\label{ex:unitbound}
\emph{Units are large.} Let $K$ be real quadratic with discriminant $d$ and fundamental unit $\varepsilon>1$. Show that
$\varepsilon\ge\frac{\sqrt{d-4}+\sqrt d}2$, and that equality holds when $d=a^2+4$ for an odd $a$. Deduce that
$R_K\ge\frac12\log d-\frac{10}{3d}$ for $d\ge5$.
\end{exercise}

\begin{beauty}
Dirichlet found the unit theorem in 1846, while listening to Easter music in the Sistine Chapel, so the story goes. Its
statement is the arithmetic twin of the finiteness of the class group: both say that $\mathcal O_K$ is ``almost'' as good as
$\Z$, one measuring the failure of unique factorisation, the other the abundance of units. The two meet in the class
number formula (Chapter~11), where only the product $h_KR_K$ appears, and they can trade size: in real quadratic fields
$h_KR_K$ is about $\sqrt d$, and nobody knows how to show that infinitely many of them have $h_K=1$ (Gauss's conjecture), which
amounts to showing that $R_K$ is often as large as $\sqrt d$. Computing $R_K$ is believed to be hard: it is one of the
problems a quantum computer solves in polynomial time (Hallgren, 2002), like factoring.
\end{beauty}

\section{Hints}

Try each exercise for a while before reading its hint, and read the hint before the solution.

\begin{hint}{ex:muK}
$\Q(\zeta_m)\subseteq K$ has degree $\varphi(m)$ (Volume~2, Exercise~10.8).
\end{hint}
\begin{hint}{ex:pell61}
$(\mathcal O/2\mathcal O)^\times$ has order $1$ or $3$; and $2\mathcal O\subseteq\Z[\sqrt d]$.
\end{hint}
\begin{hint}{ex:cyclic7}
$N(\theta)$ and $N(1+\theta)$ are read off $f(0)$ and $f(-1)$. For the roots, $2\cos\frac{4\pi}7=\theta^2-2$.
\end{hint}
\begin{hint}{ex:regcol}
$\ell(\varepsilon)=(\log\varepsilon,-\log\varepsilon)$.
\end{hint}
\begin{hint}{ex:sunits}
The kernel of the map is $\mathcal O_K^\times$. Its image has finite index in $\Z^S$, because $P^{h_K}$ is principal.
\end{hint}
\begin{hint}{ex:thue}
$N(x-y\sqrt[3]2)=x^3-2y^3$. Compute $\varepsilon^{-1}=\sqrt[3]2-1$ and $\varepsilon^{\pm2}$, $\varepsilon^{\pm3}$.
\end{hint}
\begin{hint}{ex:unitbound}
Write $\varepsilon=\frac{a+b\sqrt d}2$ with $a,b\ge1$, and use $a^2-db^2=\pm4$.
\end{hint}

\section{Solutions}

\begin{solution}{ex:muK}
\idea A root of unity of order $m$ generates a subfield of degree $\varphi(m)$, and degrees divide.
\writeup If $\zeta_m\in K$, then $\Q(\zeta_m)\subseteq K$ and $\varphi(m)=[\Q(\zeta_m):\Q]$ divides $n=[K:\Q]$ (tower law).
For $n=2$, $\varphi(m)\in\{1,2\}$ gives $m\in\{1,2,3,4,6\}$: so $\mu_K$ is $\{\pm1\}$, or contains $i$ ($K=\Q(i)$), or
contains a primitive cube or sixth root of unity ($K=\Q(\sqrt{-3})$). For $n$ odd, $\varphi(m)$ is even for $m\ge3$, so
$m\le2$ and $\mu_K=\{\pm1\}$. (Also, a field with a real embedding has $\mu_K=\{\pm1\}$, as a real root of unity is
$\pm1$.)
\end{solution}

\begin{solution}{ex:pell61}
\idea Units of $\mathcal O_K$ that are not in $\Z[\sqrt d]$ become units of $\Z[\sqrt d]$ after cubing, because the residue
ring modulo $2$ has a unit group of order at most $3$.
\writeup $N(\varepsilon)=\frac{39^2-61\cdot25}4=\frac{1521-1525}4=-1$. Let $d\equiv1\pmod4$, $\mathcal O=\Z[\omega]$,
$\omega=\frac{1+\sqrt d}2$. Then $2\mathcal O\subseteq\Z[\sqrt d]$ (as $2\omega=1+\sqrt d$), and the ring $\mathcal O/2\mathcal
O$ has $4$ elements. By Theorem~\ref{thm:kd}, it is $\mathbb F_4$ if $d\equiv5\pmod8$ ($X^2-X+\frac{1-d}4$ irreducible modulo
$2$) and $\mathbb F_2\times\mathbb F_2$ if $d\equiv1\pmod8$. Its unit group has order $3$ or $1$. So for a unit $u$, $u^3\equiv1
\pmod{2\mathcal O}$, i.e.\ $u^3\in1+2\mathcal O\subseteq\Z[\sqrt d]$.

For $d=61\equiv5\pmod8$, $\varepsilon\notin\Z[\sqrt{61}]$, and $\varepsilon^3=29\,718+3\,805\sqrt{61}$, of norm $-1$. The units
of $\Z[\sqrt{61}]$ are the $\pm\varepsilon^k$ that lie in $\Z[\sqrt{61}]$; $\varepsilon^k\in\Z[\sqrt{61}]$ iff $3\mid k$ (the class
of $\varepsilon$ in $(\mathcal O/2\mathcal O)^\times\cong\Z/3$ is non-trivial, since $\varepsilon\notin1+2\mathcal O$). Solutions of
$x^2-61y^2=1$ with $x,y>0$ are units $x+y\sqrt{61}>1$ of norm $+1$ in $\Z[\sqrt{61}]$: the $\varepsilon^{k}$ with $3\mid k$,
$k>0$, and $k$ even (as $N(\varepsilon^k)=(-1)^k$). The smallest is $\varepsilon^6=(\varepsilon^3)^2=29\,718^2+61\cdot3\,805^2+
2\cdot29\,718\cdot3\,805\sqrt{61}=1\,766\,319\,049+226\,153\,980\sqrt{61}$.
\end{solution}

\begin{solution}{ex:cyclic7}
\idea Norms of $\theta-a$ come from $f(a)$. Two units are independent when their log vectors are, which a $2\times2$ determinant
decides.
\writeup With $f=X^3+X^2-2X-1$, $N(a-\theta)=f(a)$, so $N(\theta)=-f(0)=1$ and $N(1+\theta)=-f(-1)=-(-1+1+2-1)=-1$: both are
units. The roots of $f$ are $2\cos\frac{2\pi k}7$, $k=1,2,3$ (since $\zeta+\zeta^{-1}$ for $\zeta^7=1$, $\zeta\ne1$, satisfies it,
dividing $\Phi_7(\zeta)=0$ by $\zeta^3$). By the double angle formula, $2\cos\frac{4\pi}7=\theta^2-2$, and the third root is
$-1-\theta-(\theta^2-2)=1-\theta-\theta^2$, as the roots sum to $-1$. Numerically, $\theta\approx1.2470$, $\theta_2\approx
-0.4450$, $\theta_3\approx-1.8019$. Here $r=2$ and $e_i=1$. The matrix of $\log|\cdot|$ with the third column deleted is
\[
\begin{pmatrix}\log1.2470&\log0.4450\\\log2.2470&\log0.5550\end{pmatrix}\approx\begin{pmatrix}0.2207&-0.8097\\0.8097&-0.5889
\end{pmatrix},
\]
of determinant $\approx0.2207\cdot(-0.5889)+0.8097^2\approx0.5256$. It is non-zero, so $\theta$ and $1+\theta$ are
independent. (PARI shows that they generate the whole unit group with $-1$, so $R_K\approx0.5255$.)
\end{solution}

\begin{solution}{ex:regcol}
\idea Direct computation.
\writeup $r=1$, $\varepsilon=1+\sqrt2$, $\sigma_2(\varepsilon)=1-\sqrt2=-\varepsilon^{-1}$. So $\ell(\varepsilon)=(\log\varepsilon,
-\log\varepsilon)$, and deleting either column gives $|\pm\log\varepsilon|=\log\varepsilon\approx0.8814$. The lattice
$\ell(\mathcal O^\times)=\Z(\log\varepsilon,-\log\varepsilon)$ on the line $x_1+x_2=0$ has covolume $\|\ell(\varepsilon)\|=
\sqrt2\log\varepsilon=\sqrt{r+1}\,R_K$.
\end{solution}

\begin{solution}{ex:sunits}
\idea Units are the kernel of the valuation map, and the image has full rank because powers of primes become principal.
\writeup Let $S=\{P_1,\dots,P_s\}$ and $v:\mathcal O_S^\times\to\Z^s$, $x\mapsto(v_{P_i}(x))$, a homomorphism. Its kernel is the set
of $x\in K^\times$ with $v_P(x)=0$ for every $P$, i.e.\ $x\mathcal O=\mathcal O$: the units $\mathcal O_K^\times$. Its image is a
subgroup of $\Z^s$; it contains $h_K\Z^s$, because $P_i^{h_K}=\alpha_i\mathcal O$ is principal (the class group has order $h_K$),
and $v(\alpha_i)=h_Ke_i$. So the image is free of rank $s$. The exact sequence $1\to\mathcal O_K^\times\to\mathcal O_S^\times
\to v(\mathcal O_S^\times)\to0$ splits (lift a basis of the free image), so $\mathcal O_S^\times\cong\mathcal O_K^\times\times
\Z^s\cong\mu_K\times\Z^{r+s}$ by Dirichlet.
\end{solution}

\begin{solution}{ex:thue}
\idea $x-y\sqrt[3]2$ has norm $x^3-2y^3$, so solutions are units of a special shape; all units are $\pm\varepsilon^k$, and
the sign is forced by the norm.
\writeup If $x^3-2y^3=1$, then $\alpha=x-y\sqrt[3]2$ has norm $1$ (the norm of $a+b\theta+c\theta^2$ is $a^3+2b^3+4c^3-6abc$,
Exercise~\ref{ex:normform}), so it is a unit: $\alpha=\pm\varepsilon^k$, and since $N(-1)=-1$ and $N(\varepsilon)=1$, $\alpha=
\varepsilon^k$. Conversely $\varepsilon^k=a+b\theta+c\theta^2$ with $c=0$ gives the solution $(x,y)=(a,-b)$. With
$\theta=\sqrt[3]2$ and $\theta^3=2$ one computes
\[
\begin{aligned}
\varepsilon^{-3}&=1+3\theta-3\theta^2,&\varepsilon^{-2}&=1-2\theta+\theta^2,&\varepsilon^{-1}&=-1+\theta,\\
\varepsilon&=1+\theta+\theta^2,&\varepsilon^2&=5+4\theta+3\theta^2,&\varepsilon^3&=19+15\theta+12\theta^2.
\end{aligned}
\]
(For instance $\varepsilon^{-3}=(\theta-1)^3=\theta^3-3\theta^2+3\theta-1=1+3\theta-3\theta^2$.) Among $|k|\le3$, the
coefficient of $\theta^2$ vanishes only for $k=0$, giving $(x,y)=(1,0)$, and $k=-1$, giving $(x,y)=(-1,-1)$: indeed
$(-1)^3-2(-1)^3=1$.
\end{solution}

\begin{solution}{ex:unitbound}
\idea A unit $>1$ has positive coefficients, so it is at least the smallest number of its shape allowed by the norm
equation.
\writeup Every element of $\mathcal O_K$ is $\frac{a+b\sqrt d}2$ with $a,b\in\Z$ (for $d\equiv1\pmod4$ with $a\equiv b\pmod2$;
for $d=4m$, $a+b\sqrt m=\frac{2a+b\sqrt d}2$), and a unit has $\frac{a^2-db^2}4=\pm1$. Let $\varepsilon>1$ be a unit and
$\varepsilon'=\frac{a-b\sqrt d}2=\pm\varepsilon^{-1}$, so $|\varepsilon'|<1<\varepsilon$. Then $a=\varepsilon+\varepsilon'>0$ and
$b\sqrt d=\varepsilon-\varepsilon'>0$, so $a,b\ge1$, and $a^2=db^2\pm4\ge d-4$. Hence $\varepsilon\ge\frac{\sqrt{d-4}+
\sqrt d}2$. If $d=a^2+4$ with $a$ odd, then $d\equiv5\pmod 8$ and $\frac{a+\sqrt d}2\in\mathcal O_K$ has norm $\frac{a^2-d}4=-1$:
it is a unit attaining the bound, hence the fundamental unit. Finally $\frac{\sqrt{d-4}+\sqrt d}2=\sqrt d\cdot\frac{1+\sqrt{1-4/d}}
2\ge\sqrt d\,(1-\frac2d)$, using $\sqrt{1-t}\ge1-t$ for $0\le t\le1$; and $\log(1-\frac2d)\ge-\frac2d\cdot\frac1{1-2/d}$, which is
$\ge-\frac{10}{3d}$ for $d\ge5$. So $R_K\ge\log\sqrt d-\frac{10}{3d}$.
\end{solution}

\section*{Sources and further reading}

\begin{itemize}
\item P.~G.~L.~Dirichlet, ``Zur Theorie der complexen Einheiten'', \emph{Bericht Verh.\ K\"onigl.\ Preuss.\ Akad.\ Wiss.}
(1846), 103--107. The unit theorem.
\item D.~A.~Marcus, \emph{Number Fields}, 2nd ed., Springer, 2018, Chapter~5. The proof followed here, Artin's bound, and the
units of pure cubic fields.
\item J.~Neukirch, \emph{Algebraic Number Theory}, Springer, 1999, Chapter~I, \S7. Dirichlet's theorem through the
Minkowski space.
\item J.~W.~S.~Cassels and A.~Fr\"ohlich (eds.), \emph{Algebraic Number Theory}, Academic Press, 1967 (2nd ed.\ LMS, 2010).
Artin's lemma on cubic units appears in several exercises.
\item H.~W.~Lenstra, Jr., ``Solving the Pell equation'', \emph{Notices Amer.\ Math.\ Soc.}\ 49 (2002), 182--192. The history,
the size of solutions, and the algorithms; essential reading f.
\item S.~Hallgren, ``Polynomial-time quantum algorithms for Pell's equation and the principal ideal problem'', \emph{J.\ ACM}
54 (2007). The quantum algorithm mentioned in the beauty box.
\item \emph{History.} Fermat's challenge of 1657, Brahmagupta's and Bh\=askara's work and the misattribution to Pell are
in Lenstra's article above and in A.~Weil, \emph{Number Theory: An Approach through History}, Birkh\"auser, 1984, Chapter~II.
The Sistine Chapel story is told by H.~Minkowski, ``Peter Gustav Lejeune Dirichlet und seine Bedeutung f\"ur die heutige
Mathematik'', \emph{Jahresber.\ Deutsch.\ Math.-Verein.}\ 14 (1905), 149--163; it is an anecdote, not a document.
\end{itemize}
